\documentclass[10pt,oneside]{book}
%%%% Page Info + Commands %%%%%{

%packages
\usepackage{pgfplots}
\usepackage{mathrsfs}
\pgfplotsset{compat=1.18}
\usepackage{amsmath}
\usepackage{amsfonts}
\usepackage{charter}
\usepackage{amsthm,letltxmacro}
\usepackage{amssymb}
\usepackage{fancyhdr}
\usepackage{float}
\usepackage[most]{tcolorbox}
\usepackage{enumerate}
\usepackage{enumitem}
\usepackage[dvipsnames]{xcolor}
\usepackage{changepage}
\usepackage{mdframed}
\usepackage{graphicx}

% SMALL PAGE COMMAND
\newcommand{\smallpage}{  \setlength \oddsidemargin{.25in}
            \setlength \textwidth{5.5in}
            \setlength \topmargin{0in}
            \setlength \textheight{9in}}


% Commands for sets of numbers
\newcommand{\Z}{\mathbb{Z}}\newcommand{\R}{\mathbb{R}}\newcommand{\N}{\mathbb{N}}\newcommand{\Q}{\mathbb{Q}}\newcommand{\C}{\mathbb{C}}
\newcommand{\real}[1]{\text{Re}\left(#1\right)}
\newcommand{\im}[1]{\text{Im}\left(#1\right)}

% gordie spacing and boxing commands
\newcommand{\gspc}{\vspace{0.13cm}} % 0.5 space
\newcommand{\gline}[0]{\gspc\\} % 1.5 space
\newcommand{\gbline}{\gspc\gspc\\} % double space
\newcommand{\gbox}[1]{\fbox{\parbox{\linewidth}{#1}}} % multiline box command
\newcommand{\gmod}[1]{\,(\text{mod}\,#1)}


\newcommand{\gimage}[2]{
\begin{figure}[H]
    \centering
    \includegraphics[width=#2]{#1}
\end{figure}
}

\newcommand{\gimageR}[3]{
\begin{figure}[H]
    \centering
    \includegraphics[width=#2,angle=#3]{#1}
\end{figure}
}

\newcommand{\centerit}[1]{{\begin{center} #1 \end{center}}}
\newcommand{\ul}[1]{\underline{#1}}

\newcommand{\prt}[2]{\frac{\partial#1}{\partial#2}}

%%%%%%%% command for graphics %%%%%%%%%%%%%

%}

% COLOR BOX
\newmdenv[backgroundcolor=gray!8,
  linewidth=0pt, innerleftmargin=0pt, innerrightmargin=0pt,
  skipabove=\baselineskip, skipbelow=\baselineskip
]{coloredadjust}

% PROOF
\LetLtxMacro\oldproof\proof
\let\endoldproof\endproof
\renewenvironment{proof}[1][\proofname]
  {\vspace{0.2cm}\begin{adjustwidth}{2em}{2em}
   \oldproof[#1]}
  {\endoldproof
\end{adjustwidth}}

% PROBLEM
\newenvironment{problem}[1]
  {\vspace{-0.1cm}\begin{coloredadjust}\begin{adjustwidth}{2em}{2em}
   \textit{Problem. }#1}
  {\endoldproof
\end{adjustwidth}\end{coloredadjust}\gspc}

%% DEFINITION
\newtcbtheorem[]{definition}{Definition}{enhanced,
  colback=blue!2, colframe=blue!50!black, coltitle=blue!70!black,
  fonttitle=\bfseries,
  boxrule=0pt, toprule=2pt, bottomrule=1pt, leftrule=1pt, rightrule=1pt, arc=1pt, outer arc=1pt,
  attach title to upper, title={#1},
}{dfn}

%% CRITERION
\newtcbtheorem[]{criterion}{Criterion}{enhanced,
  colback=magenta!3, colframe=magenta!60!black, coltitle=magenta!20!black,
  fonttitle=\bfseries,
  boxrule=4pt, toprule=2pt, bottomrule=1pt, leftrule=1pt, rightrule=1pt, arc=1pt, outer arc=1pt,
  attach title to upper, title={#1},
}{ctn}

%% COROLLARY
\newtcbtheorem[]{corollary}{Corollary}{enhanced,
  colback=orange!3, colframe=orange!80!black, coltitle=orange!50!black,
  fonttitle=\bfseries,
  boxrule=4pt, toprule=1pt, bottomrule=1pt, leftrule=1pt, rightrule=1pt, arc=5pt, outer arc=5pt,
  attach title to upper, title={#1},
}{ctn}

%% THEOREM
\newtcbtheorem[]{theorem}{Theorem}{enhanced,
  colback=green!2, colframe=green!40!black, coltitle=green!20!black,
  fonttitle=\bfseries,
  boxrule=4pt, toprule=3pt, bottomrule=1pt, leftrule=1pt, rightrule=1pt,arc=5pt, outer arc=5pt,
  attach title to upper, title={#1},
}{thm}




% color commands
\newcommand{\cg}[1]{\color{OliveGreen}#1\color{white}}
\newcommand{\cb}[1]{\color{BlueViolet}#1\color{white}}


\newcommand{\cpr}[1]{\left(#1\right)}
%%%%%%%%%%%%% PHYSICS SPECIFIC COMMANDS

\newcommand{\vA}{\vec{\mathbf{A}}}
\newcommand{\vB}{\vec{\mathbf{B}}}
\newcommand{\vC}{\vec{\mathbf{C}}}
\newcommand{\vZ}{\vec{\mathbf{0}}}
\newcommand{\norm}{\vec{\mathbf{n}}}
\newcommand{\pvec}[1]{\left\langle#1\right\rangle}
\newcommand{\ar}{\vec{\mathbf{r}}}
\newcommand{\arhat}{\hat{\mathbf{r}}}
\newcommand{\vv}{\vec{\mathbf{v}}}
\newcommand{\delfull}{\pvec{\prt{}{x}, \prt{}{y},\prt{}{z}}}
\newcommand{\delinput}[3]{\pvec{\prt{}{x}\cpr{#1}, \prt{}{y}\cpr{#2},\prt{}{z}\cpr{#3}}}
\newcommand{\dps}{\displaystyle}
\newcommand{\delmatrix}[3]{\begin{bmatrix}\dps\hat x&\dps\hat y&\dps\hat z\\\dps\prt{}x&\dps\prt{}y&\dps\prt{}z\\\dps#1&\dps#2&\dps#3\end{bmatrix}}

%%%%%%%%%%%%%%%%%%%%%%%%%%%%
%%%%%%%%%%%%%%%%%%%%%%%%%%%%%%
%%%%%%%%%%%%%%%%%%%%%%%%%%%%%
%%%%%%%%%%%%%%%%%%%%%%%%%%%%%
%%%%%%%%%%%%%%%%%%%%%%%%%%%%%%
%%%%%%%%% begin doc

%%%% Page Setup %%%%%%%
\smallpage\sloppy
\pagestyle{fancy}
\lhead{\large{\textbf{HW2 - Theory of Computation}}}
\setlength{\headheight}{10pt}



\begin{document}

\begin{enumerate}
  \item Write out the formal definition for the NFA with state diagram pictured below:
  \gimage{image.png}{5cm}
  \begin{problem}
    Let $A = \{0\}$. Let $B = \{1\}$. This NFA can be described as:
    \begin{align*}
      A^*\cup(A \circ B^*)
    \end{align*}
    We write out the formal classification of this NFA as:
    \begin{align*}
      \Sigma = \{0, 1\}, \quad q_0 = q_1, \quad F = \{q_2,q_3\},\quad Q = \{q_1, q_2, q_3\}
    \end{align*}
    \begin{center}
      \begin{tabular}{c||c|c|c}
        $\delta$ & $q_1$ & $q_2$ & $q_3$\\\hline\hline
        0 & $q_1$ & $q_2$ & $\emptyset$\\\hline
        1 & $\emptyset$ & $\emptyset$ & $q_3$ \\\hline
        $\varepsilon$ & $q_2$ & $\emptyset$ & $\emptyset$
      \end{tabular}
    \end{center}
  \end{problem}
  \item Draw a state diagram for NFAs that recognizes the fllowing set of binary strings
  \begin{enumerate}
    \item The language of strings that contains exactly one 1 followed by an odd number of zeros. 
    \begin{problem}
      I'll assume that we can have any number of zeros beforehand because it doesn't specify that in the question.
      \gimage{image2.png}{7cm}
    \end{problem}
    \item Binary strings that contain at least one 1 and where all 0s are followed by a 1.
    \begin{problem}
      Here, we just ensure that all 0s are followed by a 1, which puts us into the accepting state.
      \gimage{image3.png}{6cm}
    \end{problem}
  \end{enumerate}
  \item Consider the NFA with the previous state diagram. 
  \gimage{image4.png}{7cm}
  \begin{enumerate}
    \item Give five examples of strings recognized and five examples of strings not recognized by this NFA.
    \begin{problem}
      Here are five strings that are recognized:
      \begin{align*}
        \{00, 11, 0011, 0000, 1111\}
      \end{align*}
      Here are five strings that are not recognized:
      \begin{align*}
        \{01, 10, 010, 101, 0010\}
      \end{align*}
    \end{problem}
    \item Describe the language recognized by this NFA.
    \begin{problem}
      It recognizes all nonempty binary strings where every 0 or 1 is followed by a 0 or 1, respectively. 
    \end{problem}
  \end{enumerate}
  \item Find DFAs that recognizes the same language as the NFAs with diagrams below. What is their language?
  \begin{enumerate}
    \item {$\cdots$}\\
    \gimage{image5.png}{5cm}\vspace{-0.5cm}
    \begin{problem}
      Here is the DFA that I envisioned:
      \gimage{image6.png}{5cm}
      It describes the language of all binary strings that do not start with zero. 
    \end{problem}
    \item $\cdots$\\
    \gimage{image7.png}{5cm}
    \begin{problem}
      Here is the DFA that I envisioned for this problem:
      \gimage{image8.png}{5cm}
      It describes all binary strings where no 1s follow a 0. 
    \end{problem}
  \end{enumerate}
  \item Consider the languages $A=\{0^k:k\ge 0\}$ (strings with zero or more 0s, but no 1s) and $B = \{01\}$ (only one string, 01).
  \begin{enumerate}
    \item Draw an NFA that recognizes $A$. Draw an NFA that recognizes $B$. 
    \gimage{image9.png}{5cm}
    \item Draw the diagram for an NFA that recognizes $A^c$.
    \gimage{image10.png}{5cm}
    \pagebreak
    \item Draw the diagram from an NFA that recognizes the language $A\cup B$. 
    \gimage{image11.png}{5cm}
    \item Draw the diagram for an NFA that recognizes the language $A\circ B$
    \gimage{image12.png}{5cm}
    \item Draw the diagram for an NFA that reocgnizes the language $B^*$. 
    \gimage{image13.png}{5cm}
  \end{enumerate}
\end{enumerate}


\end{document}